\chapter{CONCLUSION AND FUTURE WORK}
\label{chap:conclusion}


%─────────────────────────────────────────────────────────────────────────────
\section{CONCLUSION}
%─────────────────────────────────────────────────────────────────────────────

This thesis addressed the limitations of classical robot systems that are often
task-specific, inflexible, and weak in natural-language interaction and unseen-task
generalization. To solve these issues, the work proposed a Vision-Language-Action
(VLA) framework that integrates three modules: a VLM core, a Vision Expert, and an
Action Expert.

The project achieved the current-stage objectives:
\begin{enumerate}[label=\arabic*)]
    \item Developed an end-to-end VLA system for robot manipulation.
    \item Integrated VLM + Vision Expert + Action Expert into one coherent pipeline.
    \item Enabled robots to follow natural-language instructions.
    \item Improved generalization performance on unseen tasks.
    \item Evaluated the system on representative manipulation benchmarks.
\end{enumerate}

Experimental results showed that explicit module integration improved action validity
and task completion compared with partial baselines. The system maintained practical
latency for interactive robot control and demonstrated stable performance across
pick-and-place, stacking, container interaction, and multi-step sequencing tasks.

Overall, the results support the central claim of this thesis: a carefully integrated
VLA architecture can provide a practical path toward general-purpose intelligent robot
manipulation under real-world compute and data constraints.


%─────────────────────────────────────────────────────────────────────────────
\section{LIMITATIONS}
%─────────────────────────────────────────────────────────────────────────────

Despite strong overall performance, several limitations remain.

\begin{itemize}
    \item Generalization to highly novel long-horizon tasks still degrades compared
    with seen-task performance.
    \item Data scale and diversity remain limited relative to real-world task
    complexity.
    \item Action sequencing errors can appear when instructions are ambiguous or
    underspecified.
\end{itemize}

These limitations indicate that stronger temporal reasoning and broader data coverage
are necessary for robust deployment in unconstrained environments.


%─────────────────────────────────────────────────────────────────────────────
\section{FUTURE WORK}
%─────────────────────────────────────────────────────────────────────────────

Future development will focus on the following directions:

\begin{enumerate}[label=\arabic*)]
    \item Expand robot demonstration data with broader object, scene, and instruction
    diversity to strengthen unseen-task generalization.
    \item Improve long-horizon planning in the Action Expert for more reliable
    multi-step task execution.
    \item Optimize model efficiency further for lower-latency operation on embedded
    hardware.
    \item Extend evaluation to more complex real-world manipulation scenarios with
    stricter robustness requirements.
    \item Introduce stronger safety checks and action-validation layers before physical
    execution.
\end{enumerate}

These directions will support the transition from a research prototype to a more
widely deployable general-purpose robotic assistant.
